The notation below is used consistently throughout the paper and the
records. Implementations may differ in scale conventions (e.g.\
network-total vs.\ per-neuron current); where a number is quoted, the
paper states the convention.

\begin{center}
\small
\begin{tabular}{ll}
\toprule
Symbol & Meaning \\
\midrule
\(N = 52\) & internal/output neurons (ids 24--75) \\
\(a = 52\) & synaptic amplitude (current per spike \(= a\cdot w\)) \\
\(w\) & synaptic weight, \([0,1]\) \\
\(t_e = 0.8\) & M2 per-neuron excitatory budget \\
\(P_i\) & per-neuron protected (consolidated) mass \\
\(p_{\max} = 0.75\) & protected cap fraction; cap \(p_{\max} t_e = 0.60\) \\
\(W_i = t_e - P_i\) & working (unconsolidated) budget \\
\(R_i = I_{p,i}/(I_{p,i}+I_{w,i})\) & protected fraction of delivered input current (allocation gate) \\
\(I_{W,i}\) & working (unconsolidated) input-channel current delivered to \(i\) \\
\(I_{\mathrm{rec}}, I_{\mathrm{inh}}\) & recurrent excitatory / M6 inhibitory current \\
\(v_i, v_{\mathrm{th}}=1, \tau_m=20\) & membrane potential, threshold, membrane time constant (ms) \\
\(u_i\) & V2.1 slow state; \(+\beta\)/spike, decay \(\tau_s\) \\
\(\beta = 0.0046875, \tau_s = 5000\) & the e24 cell used in the CLLA era \\
\(k_g = 8.0\) & frozen recruitment-gain magnitude \\
\(i_{\mathrm{boost},i}\) & recruitment-gain current (registration A; capped in B) \\
\(\theta_{\mathrm{perm}} = 0.05, w_{\mathrm{c,perm}} = 0.02\) & M3 permanence threshold / bootstrap weight \\
\end{tabular}
\end{center}

Conventions: ``posts/n'' is the fraction of the 52 neurons that fired
at least once in a presentation window; ``burst rate'' is the
presentation-window mean internal spike rate; ``protected mean'' is the
mean weight over live consolidated synapses of a cohort; ``drive end''
is the snapshot at \(t=84{,}000\) ms (the e24 convention);
per-presentation timing is 500 ms on / 1500 ms off. Where the records
report network-total currents (e.g.\ 81.38 total vs.\ 1.565 per neuron
at first C), the paper quotes the per-neuron value.